% Page 1
\begin{ApplicationFormPage}
  \ApplicationPart{79.7}{第一部分\quad 国内外现状及趋势分析}
  \ApplicationFrame{103.2002}{762.9986}{103.4387,176.5911,744.6919,763.2410}
  \ApplicationLabel{106.1}{%
    一、国外研究现状及趋势（包括该方面国外的总体研究水平、最新进展和\\
    发展预期，国外从事该方面研究的主要机构及其研究方法、成果；相关研\\
    究的学术影响、知识产权申请和授予、技术标准制定及产业应用状况，并\\
    列出主要的文献，专利、标准名称。限 1500 字以内。）}
  \ApplicationAnswer{180.0}{560.5}{\ForeignResearchContent}
  \ApplicationLabel{747.1}{二、国内研究现状及趋势（包括国内从事该方面研究的主要机构及其研究}
\end{ApplicationFormPage}

% Page 2
\begin{ApplicationFormPage}
  \ApplicationFrame{71.9995}{722.9515}{72.2419,109.0391,723.1901}
  \ApplicationLabel{75.1}{%
    方法、研究水平；国内外研究方法、研究进展的对比分析；该方面研究可\\
    能取得的突破及预期成果。限 1500 字以内。）}
  \ApplicationAnswer{113.0}{606.0}{\DomesticResearchContent}
\end{ApplicationFormPage}

% Page 3
\begin{ApplicationFormPage}
  \ApplicationPart{79.7}{第二部分\quad 研究内容及目标}
  \ApplicationFrame{103.2002}{749.1503}{103.4387,131.7884,423.1893,749.3888}
  \ApplicationLabel{112.0}{一、项目目标及考核指标}
  \ApplicationLabel{140.7}{%
    （一）申报项目与所属指南方向的关联关系（包括项目与所属指南方向的\\[13.2bp]
    匹配性，对指南方向目标的支撑作用。限 500 字以内。）}
  \ApplicationAnswer{195.0}{224.0}{\GuideAlignmentContent}
  \ApplicationLabel{432.0}{（二）项目目标及考核指标、评测方式/方法（限 800 字以内）}
  \ApplicationAnswer{460.0}{285.0}{\ObjectivesContent}
\end{ApplicationFormPage}

% Page 4
\begin{ApplicationFormPage}
  \ApplicationFrame{71.9995}{749.0016}{72.2419,282.1917,310.5414,749.2402}
  \ApplicationLabel{81.3}{%
    （三）项目成果的呈现形式及描述（重点是新产品、新设备、新品种、新\\[13.2bp]
    标准、新工艺，包括专利、标准、成果等，限 500 字以内）}
  \ApplicationAnswer{135.0}{143.0}{\OutcomeFormsContent}
  \ApplicationLabel{290.7}{二、项目研究内容、研究方法及技术路线}
  \ApplicationLabel{319.3}{（一）项目的主要研究内容（限 1000 字以内）}
  \ApplicationAnswer{347.0}{398.0}{\ResearchContent}
\end{ApplicationFormPage}

% Page 5
\begin{ApplicationFormPage}
  \ApplicationFrame{71.9995}{758.7983}{72.2419,759.0408}
  \ApplicationLabel{81.3}{（二）项目拟采取的研究方法}
  \ApplicationLabel{112.0}{%
    \ApplicationExactLine{453.49}{1、项目研究拟解决的问题，及拟采用的方法、原理、机理、算法、模型等}\\[13.2bp]
    （研究方法、技术路线可用图表描述。限 800 字以内）}
  \ApplicationAnswer{166.0}{498.0}{\ResearchMethodsContent}
  \ApplicationLabel{674.0}{2、项目研究方法（技术路线）的可行性、先进性分析（限 500 字以内）}
  \ApplicationAnswer{705.0}{50.0}{\MethodFeasibilityContentA}
\end{ApplicationFormPage}

% Page 6
\begin{ApplicationFormPage}
  \ApplicationFrame{71.9995}{568.7509}{72.2419,540.7414,568.9894}
  \ApplicationAnswer{78.0}{458.0}{\MethodFeasibilityContentB}
  \ApplicationLabel{549.3}{三、主要预期创新点}
\end{ApplicationFormPage}

% Page 7
\begin{ApplicationFormPage}
  \ApplicationFrame{71.9995}{713.9018}{72.2419,402.1881,430.5417,714.1404}
  \ApplicationLabel{81.3}{%
    （围绕基础前沿、共性关键技术或应用试验等层面，简述项目预期的主要\\[13.2bp]
    创新点。具体内容应包括该项创新的基本形态及其前沿性、时效性等，并\\[13.2bp]
    说明是否具备方法、理论和知识产权特征。每项创新点的描述限 200 字以\\[13.2bp]
    内）}
  \ApplicationAnswer{205.0}{193.0}{\InnovationContent}
  \ApplicationLabel{410.7}{四、预期经济社会效益}
  \ApplicationLabel{439.3}{项目的科学、技术、产业预期指标及社会、经济效益（限 500 字以内）}
  \ApplicationAnswer{467.0}{243.0}{\BenefitsContent}
\end{ApplicationFormPage}

% Page 8
\begin{ApplicationFormPage}
  \ApplicationPart{79.7}{第三部分\quad 申报单位及参与单位研究基础}
  \ApplicationFrame{103.2002}{742.0990}{103.4387,131.7884,742.3415}
  \ApplicationLabel{112.0}{一、申报单位的已有工作基础、研究成果、研究队伍等}
  \ApplicationLabel{140.7}{%
    （一）申报单位在该研究方向的前期任务承担情况、相关研究成果（应说\\[13.2bp]
    明前期任务取得的主要成果、指标及效果，限 800 字以内）}
  \ApplicationAnswer{195.0}{543.0}{\PriorWorkContent}
\end{ApplicationFormPage}

% Page 9
\begin{ApplicationFormPage}
  \ApplicationFrame{71.9995}{749.0486}{72.2419,749.2910}
  \ApplicationLabel{81.3}{%
    \ApplicationExactLine{460.57}{（二）负责人及项目主要骨干人员的科研水平及主要成果（包括工作简历、}\\[13.2bp]
    主要学术业绩、近五年主持的与申请项目相关的各类国家科技计划项目情\\[13.2bp]
    况，奖励、论文、专利等重点成果取得情况，限 1200 字以内）}
  \ApplicationAnswer{175.0}{570.0}{\TeamAchievementsContent}
\end{ApplicationFormPage}

% Page 10
\begin{ApplicationFormPage}
  \ApplicationFrame{71.9995}{749.0016}{72.2419,375.6413,453.7408,749.2402}
  \ApplicationLabel{81.3}{%
    （三）申报单位相关科研条件支撑状况（包括实验平台、大型仪器设备、\\[13.2bp]
    国家（重点）实验室、国家工程（技术）中心、重大科学工程等研究基地\\[13.2bp]
    情况，限 500 字以内）}
  \ApplicationAnswer{175.0}{196.0}{\ResearchConditionsContent}
  \ApplicationClippedLabel{384.7}{375.6413}{453.7408}{%
    二、协作单位的选择原因及其优势（企业单位重点说明与项目相关的技术\\[13.2bp]
    开发与应用水平、管理能力和经济实力，与本项目相关产品开发、生产、\\[13.2bp]
    经营情况等；非企业类单位重点说明在项目相关技术方面的研发能力、技}
  \ApplicationAnswer{460.0}{285.0}{\PartnerAdvantagesContent}
\end{ApplicationFormPage}

% Page 11
\begin{ApplicationFormPage}
  \ApplicationFrame{71.9995}{742.0990}{72.2419,100.5917,508.8915,537.2412,742.3415}
  \ApplicationLabel{81.3}{三、申请单位与协作单位的任务分工}
  \ApplicationAnswer{105.0}{400.0}{\TaskDivisionContent}
  \ApplicationLabel{518.0}{四、相关的国际合作与交流（限 300 字以内）}
  \ApplicationAnswer{542.0}{196.0}{\InternationalCooperationContent}
  \ApplicationPart{749.5}{第四部分\quad 进度安排}
\end{ApplicationFormPage}

% Page 12
\begin{ApplicationFormPage}
  \ApplicationFrame{71.9995}{737.1518}{72.2419,135.1400,737.3903}
  \ApplicationLabel{81.3}{%
    \ApplicationExactLine{460.57}{一、进度安排（包括项目主要研究任务的研发进度、年度及重点节点安排、}\\[13.2bp]
    中期目标等。可用图表描述，限 1000 字以内。）}
  \ApplicationAnswer{141.0}{592.0}{\ScheduleContent}
\end{ApplicationFormPage}

% Page 13
\begin{ApplicationFormPage}
  \ApplicationPart{79.7}{第五部分\quad 项目组织实施、保障措施及风险分析}
  \ApplicationFrame{103.2002}{761.2504}{103.4387,142.9891,446.5409,509.4390,761.4890}
  \ApplicationClippedLabel{112.0}{103.4387}{142.9891}{%
    \ApplicationExactLine{460.57}{一、项目组织实施机制（包括项目及分任务（课题）的内部组织管理方式、}\\[13.2bp]
    协调机制等，限 800 字以内）}
  \ApplicationAnswer{165.0}{277.0}{\OrganizationContent}
  \ApplicationLabel{455.3}{%
    二、保障措施（项目实施的外部政策、组织和资源支撑条件，限 500 字以\\[13.2bp]
    内）}
  \ApplicationAnswer{515.0}{242.0}{\SafeguardsContent}
\end{ApplicationFormPage}

% Page 14
\begin{ApplicationFormPage}
  \ApplicationFrame{71.9995}{760.4995}{72.2419,100.5917,360.4398,423.3418,760.7381}
  \ApplicationLabel{81.3}{三、知识产权对策、成果管理及合作权益分配（限 500 字以内）}
  \ApplicationAnswer{106.0}{250.0}{\IntellectualPropertyContent}
  \ApplicationLabel{369.3}{%
    四、风险分析及对策（包括潜在的技术风险、市场风险、政策风险，实施\\[13.2bp]
    过程中的制约因素等，以及相应对策措施，限 500 字以内）}
  \ApplicationAnswer{428.0}{328.0}{\RiskContent}
\end{ApplicationFormPage}
